\begin{description}
\item[(1)] Total area within the ellipse (actually, half-light ellipse) of
  NGC 1052-DF2 is:
  \begin{equation*}
  A_{DF2} = \pi a b = 0.85\pi \times 22.6''^{\,2}
  = 1363.9~\mathrm{arcsec^2}
  \tag*{(2 points)}
  \end{equation*}
  Magnitude of the part of the galaxy within the ellipse:
  \begin{equation*}
  m_{ell} = SB - 2.5 \times \log(A_{DF2})
  = 24.7 - 2.5 \times \log(1363.9) = 16.9
  \tag*{(2 points)}
  \end{equation*}
  Total magnitude of the whole galaxy:
  \begin{equation*}
  m_{gal} = m_{ell} - 2.5 \times \log(2) = 16.1
  \tag*{(2 points)}
  \end{equation*}
\item[(2)] Absolute mag of the galaxy:
  \begin{equation*}
  M_{0,gal} = m_{gal} + 5 - 5 \times \log(D) = -15.4
  \tag*{(2 points)}
  \end{equation*}
  Convert to solar luminosity:
  \begin{equation*}
  \frac{L_{gal}}{L_\odot} = 10^{-0.4(M_{0,gal} - M_{0,\odot})}
  = 1.2 \times 10^{8}
  \tag*{(2 points)}
  \end{equation*}
  Thus, the total stellar mass should be:
  \begin{equation*}
  \frac{M_{gal}}{M_\odot} = \frac{L_{gal}}{L_\odot} \times \frac{M}{L}
  = 2.4 \times 10^{8}
  \tag*{(2 points)}
  \end{equation*}
\item[(3)] Mean galactocentric distance of globular clusters:
  \begin{equation*}
  r_{gc} = \frac{\theta D}{206265} = 7.6~\mathrm{kpc}
  \tag*{(2 points)}
  \end{equation*}
  Using viral theorem: % sic: "viral" for "virial" as in the original
  \begin{gather*}
  \langle 2K \rangle + \langle U \rangle = 0, \\
  \langle K \rangle = \tfrac{1}{2}M\langle v^2 \rangle
  = \tfrac{1}{2}M\sigma^2, \\
  \langle U \rangle = \frac{3}{5}\,\frac{GM^2}{r}
  \tag*{(4 points)}
  \end{gather*}
  Thus the maximal dynamical mass of this system should be:
  \begin{equation*}
  M_{dyn} = \frac{5 r_{gc}\sigma^2}{3G} = 2.1 \times 10^{8}\,M_\odot,
  \tag*{(3 points)}
  \end{equation*}
  even less than the total stellar mass within this radius.
\item[(4)] From Questions 2 and 3, we know:
  \[
  M_{star} \propto L_{gal} \propto F_{gal} \times D^2
  \]
  However,
  \begin{equation*}
  M_{dym} \propto D % sic: "dym" as in the original
  \tag*{(2 points)}
  \end{equation*}
  Thus
  \begin{equation*}
  M_{dym}/M_{star} \propto 1/D
  \tag*{(2 points)}
  \end{equation*}
  If the distance measured is smaller by a certain factor, the stellar mass to
  dynamical mass ratio would increase by the same factor, thus the dark matter
  in NGC 1052-DF2 would not be as deficient as Dragonfly team claimed.
\end{description}
